% !PRINTED 304
\weberpage{304}

\begin{center}
\S.\,100.

\gesperrt{Die Newton'sche Regel.}
\end{center}

Eine Ergänzung der vorstehenden Regel, die eine weitere Annäherung an
die genaue Festsetzung der Grenzen der Wurzeln gestattet, giebt die
jetzt zu besprechende, von \gesperrt{Newton} herrührende, aber erst von
\gesperrt{Sylvester} bewiesene Vorschrift\footnote{\gesperrt{Newton},
Arithmetica universalis. \gesperrt{Sylvester}, Transactions of the
R. Irish Academy, t.\,24 (1871). Phil. Mag., 4.\,ser., t.\,31. Vgl. auch
\gesperrt{Petersen}, Theorie der algebr. Gleichungen. Kopenhagen 1878.}.

Es sei wieder
\weq{1}{f(x) = a_0 x^n + a_1 x^{n-1} + \cdots + a_{n-1} x + a_n = 0}
die zu untersuchende Gleichung. Wir setzen, indem wir mit $\Pi(m)$ wie
immer das Product $1.\,2.\,3 \ldots m$ und mit $f^{(\nu)}(x)$ die
$\nu^{\text{te}}$ Derivirte bezeichnen,
\weq{2}{f_\nu(x) = \frac{\Pi(n - \nu)}{\Pi(n)} f^{(\nu)}(x),}
\weq{3}{F_\nu(x) = f_\nu^2(x) - f_{\nu+1}(x) f_{\nu-1}(x),}
mit dem Zusatz, dass
\weq{4}{F_0 = F = 1, \quad F_n = f_n^2,}
also, worauf es wesentlich ankommt, positiv sein sollen.

Für die Abgeleiteten der Functionen $f_\nu$, $F_\nu$ erhalten wir aus
(2) und (3)
\weq{5}{f_\nu'(x) = (n - \nu) f_{\nu+1},}
\weq{6}{f_\nu F_\nu'(x) = (n - \nu - 1)(F_\nu f_{\nu+1} + F_{\nu+1} f_{\nu-1}).}

Da es auf die positiven Zahlenfactoren nicht ankommt, so können wir die
Budan'sche Reihe durch die Reihe der Functionen
\[
f, f_1, f_2 \ldots f_n
\]
ersetzen.

Die \gesperrt{Newton'sche} Regel macht nun von der Doppelreihe Gebrauch:
\weq{7}{
\begin{array}{l}
f,\ f_1,\ f_2 \ldots f_n \\
F,\ F_1,\ F_2 \ldots F_n.
\end{array}
}
